\documentclass[11pt]{article}
\usepackage{amsmath,amssymb,amsthm,geometry}\geometry{margin=1in}
\newtheorem{definition}{Definition}\newtheorem{theorem}[definition]{Theorem}\newtheorem{proposition}[definition]{Proposition}
\title{Family contexts, polynomial content and directed crossing coherence, v3}
\date{16 November 2026}
\begin{document}\maketitle
\paragraph{Audit and exact claim.} Relative-assembly theorem. This separate version additionally spells out the corner functors of the higher diagrams; v2 SHA256 \texttt{91ea556540467ee613ff5c33fa1136d4c5be65a105d2180a85d32e9c46fa3971} remains unchanged. It corrects the carrier-functor nonfaithfulness sentence in v1 SHA256 \texttt{1a3174e73fd7b683446ca9110fd5524d61495604d5ea81e86a736a734a8d31f6}, which remains unchanged, and makes the strong context-module source restriction explicit. The source specialization is the normalized product-rig context and constant ordered-polynomial content of \emph{Globally Graded Open Interaction}, SHA256 \texttt{182a8462b3a7732967ad59f774dc2dd91090bf60c09a3329c9dd8c31e5bbd07b}. The reviewed initializer-family skew composite is \texttt{c91ba18085731d30c7b735e726062fc1756733e0065f6b8df4420c0b54565090}. This manuscript specifies an enlarged context module, an unchanged independently full content module, an all-coupling crossing, the directions of both multiplication comparisons, and concrete generating coherence diagrams. It does not assert an unnamed general-purpose distributive-law theorem or paid-system composition. All raw carriers are finite; full shape libraries are small in an adequate universe. Conditional plans at zero prefixes are not separate data.

\paragraph{Categorical dimension.} $M_{\mathsf F},N,H_{\mathsf F}$ are category-valued relative hom constructions. The symbol $\odot$ between modules denotes the same balanced relative tensor/native coend as in the source, while the mixed binary product is specified separately by its normal-form formula. The crossing is a functor on those balanced hom categories. The unit and multiplication comparisons below are natural transformations between the indicated composite functors, with invertible components only where stated. The higher assertions are equations between pasted natural transformations, verified on their complete native components and their naturality. No equivalence between this specified data and every axiom of an independently named skew distributive-law formalism is claimed without an additional matching theorem.

\section{The context and content modules independently}
An object of the context category $\mathcal K_A^{\mathsf F}$ is, at every input $a$, a finite common raw carrier $W_a$, augmented context objects $L_{a,w}$, a nonempty convex family $\mathcal A_a\subseteq\operatorname{Prob}(W_a)$, and a distinguished member. A context arrow has common raw functions and actual augmented coefficient arrows at all atoms; it pushes the entire source family into the target family and preserves the distinguished law. Define $e_AP=\coprod_aW_a$. The context category retains distinct actual coefficient arrows and distinct law-family objects. The carrier functor forgets those fields and may identify them; no faithfulness of $e_A$ is assumed.

Put
\[
 M_{\mathsf F}(A,B)=(e_A\downarrow B).
\]
A context-process body assigns one output label to each raw atom. Its dependent accumulation uses the full pair carrier $(w,z)$ and coefficient $L_w\otimes K_{w,z}$. Its law family consists of all laws whose outer marginal is admitted and whose conditional inner law at each positive outer atom is in the specified inner family. Inner families and raw carriers may depend on the full declared outer raw atom. At a zero outer atom there is no chosen conditional kernel; every raw pair and coefficient still exists. The distinguished law is dependent independent accumulation of distinguished members.

This accumulation is functorial: merged outer atoms give convex mixtures of pushed inner laws with compatible input types. It is strongly associative on dependent context paths. Both bracketings retain the entire path $(w,z,t)$; the last family is indexed by the same full prefix $(w,z)$ on both sides, not by an observer forgetting part of that prefix. Thus their law tests coincide. Actual coefficients use tensor associators. The singleton context unit is strong, and pure reindexing and middle balance retain the same raw pairs. In particular the unary module $M_{\mathsf F}$ does not exhibit the sibling-conditioning obstruction.

Independently retain the source content category of entire weighted ordered collections $C=(C_n,(R_c)_{c\in C_n})$, with extension $F_CB=\coprod_nC_n\times B^n$, full arity-indexed shape maps and actual coefficient maps on all shapes. Put
\[
 N(A,B)=(A\downarrow F_{(-)}B).
\]
Its multiplication is the original full substitution $C\circ D$, with the original addressed shape bijections and actual coefficient associators; its unit is the unary unit collection. Its pure actions and balance are unchanged. Native content maps preserve ordered positions, so they do not contract or duplicate context occurrences. A broader polynomial map class is not covered.

The mixed normal forms are
\[
 H_{\mathsf F}(A,B)=(e_A\downarrow F_{(-)}B).
\]
They retain one full collection, a family context object and a body on every raw atom. In the same native relative-coend convention as the source, $M_{\mathsf F}\odot N$ has these normal forms: reduce a representative with intermediate set $S$, context body $g:e_AP\to S$ and content body $h:S\to F_CB$ to $(P,C,hg)$. Conversely use intermediate set $e_AP$ and identity context body. Pure middle balance gives the two inverse normal-form comparisons. Native arrows retain the full context arrow and whole grade arrow throughout this reduction; equality of carrier functions alone is not imposed on them. Thus the old normal-form argument does not rely on faithfulness of $e_A$.

\section{The all-coupling crossing}
For $N$-data $f(a)=(c_a,b_1,\ldots,b_n)$ and family context data at each $b$, with raw $W_b$, coefficients $L_{b,w}$, families $\mathcal A_b$ and output body $g(b,w)$, define a mixed output with raw
\[
 \prod_{i=1}^nW_{b_i},\qquad
 \mathcal C_a=\{\lambda:\text{each coordinate marginal lies in }\mathcal A_{b_i}\},
 \qquad \bigotimes_i L_{b_i,w_i}.
\]
For $n=0$ this is the singleton unit row. The body is $(c_a,g(b_1,w_1),\ldots,g(b_n,w_n))$, and the whole content collection is the original $C$, not its realized subcollection. The distinguished member is the product of distinguished marginals. This gives the crossing functor, written
\[
 \chi:N\odot M_{\mathsf F}\longrightarrow M_{\mathsf F}\odot N
\]
using the stated normal forms on the target.

Its action on native maps is the full coordinate raw map, tensor of actual context arrows and original entire content arrow. Pushforward sends every admitted coupling to an admitted coupling. A noninjective output/pure-middle function may identify labels but never merges the distinct occurrence positions. Both representatives of pure middle balance therefore have the same factors and the same complete law family. Thus the crossing descends through the native balance relations and commutes with the source pure input/output actions. No universal selected coupling or selected off-image extension is invoked.

\section{Units and the two different multiplication strengths}
\paragraph{Context unit.} Crossing a content operation past singleton context units gives a forced singleton raw tuple and unit tensor factors. Deleting them is an invertible native comparison, identity on the entire content collection.

\paragraph{Content unit.} Crossing the unary unit content past a family context has one coordinate; its coupling family is precisely the original law family. The raw and actual coefficient maps are unit isomorphisms. These two unit comparisons retain zero raw atoms of nonunit context objects; only forced singleton unit coordinates are deleted.

\begin{proposition}[Content-multiplication comparison is invertible]
Crossing a composite of two content operations past a family context agrees, by an actual native isomorphism, with first crossing the inner content and then the outer content, with their strong multiplication comparisons.
\end{proposition}
\begin{proof}
At an input, deterministic content data specify a finite tree of leaf context occurrences $(i,j)$. The direct crossing admits every law on the complete leaf tuple whose leaf marginals belong to their assigned families. The iterated crossing first assigns each $i$ block its family of all such within-block couplings, then admits any coupling of those block tuples whose block marginal belongs to that family. These are exactly the same conditions: a block law belongs iff its leaf marginals are admitted. Restriction of the entire flat law and retention of its complete outer coupling are inverse constructions. This works on all tuples including zeros; no selected conditional on a zero block is reconstructed. Context coefficients are the same complete tensor, canonically regrouped. The whole content collection is $C\circ D$ on all formal shapes and has the original substitution comparison, including every unused coefficient arrow.\end{proof}

\begin{proposition}[Context-multiplication comparison is directed]
For a content operation followed by two family-context stages, the route crossing the content successively past the two stages maps natively into the route first composing those context stages and then crossing once. In schematic notation its direction is
\[
 (\mu_M N)(M_{\mathsf F}\chi)(\chi M_{\mathsf F})
 \ \Longrightarrow\ \chi(N\mu_M),                         \tag{1}
\]
with the existing strong reassociations understood. This comparison need not be invertible.
\end{proposition}
\begin{proof}
Fix the ordered content occurrences. Let $x_i$ be their first-context witnesses and $y_i$ their second-context witnesses; the type of $y_i$ depends on the output of $x_i$. The successively crossed route admits laws on the entire tuple $(\mathbf x,\mathbf y)$ such that the $\mathbf x$ marginal has admitted coordinate marginals and, at each positive whole $\mathbf x$, each $y_i$ conditional marginal belongs to its prescribed second family. The once-crossed route instead asks that each local pair $(x_i,y_i)$ marginal be an admitted unary context composite: its $x_i$ marginal is admitted and its conditional $y_i$ law at positive $x_i$ is admitted. The first test implies the second by convexly mixing over the other first-context coordinates. These mixtures have the same local input label because $x_i$ is retained. This proves admission of the identity full-tuple map. Its actual context coefficient map is the labeled tensor reassociation/permutation from stage order to occurrence order. The whole grade is the same content $C$, with identity full grade map.

For strictness use binary content and fair singleton first/second context rows. The law with independent fair $x_1,x_2$ and $y_1=x_2$, $y_2=x_1$ has both local pair marginals independent fair, but has deterministic $y_i$ conditional on the whole first tuple. It belongs only to the once-crossed family. All sixteen raw tuples remain, twelve with zero mass. Thus no inverse with the specified raw rebracketing exists.\end{proof}

This distinction is not an artifact of tensor coefficients: their displayed arrows are isomorphisms. It is the distinction between conditioning on a whole sibling tuple and on the relevant local prefix. Nor does (1) forget a separately supplied zero-branch kernel: no such kernel is a field of these families.

\section{Named generating coherence diagrams}
The following is the coherence claim proved here. Its generators are pure balance, the strong multiplications and units of $M_{\mathsf F}$ and $N$, the crossing, the two crossing unit comparisons, the invertible content-multiplication comparison, and directed comparison (1). Every comparison is an actual common-raw map with its specified full coefficient/grade map. The claim is not a substitute for checking that an external formal definition uses precisely these generators and orientations.

More explicitly, the three-content comparison is a diagram of functors out of $N\odot N\odot N\odot M_{\mathsf F}$; the three-context comparison is a diagram of natural transformations between the fully staged crossing functor and the single crossing after triple context multiplication, out of $N\odot M_{\mathsf F}\odot M_{\mathsf F}\odot M_{\mathsf F}$. The mixed comparison is out of $N\odot N\odot M_{\mathsf F}\odot M_{\mathsf F}$, comparing content multiplication followed by the two-stage-to-one-stage context comparison with first applying the content-crossing multiplication comparison and then that directed comparison inside and outside the retained blocks. Strong relative reassociations place each pair of routes over the same indicated domain and target functors. The two unit diagrams are obtained from the words $N\odot I_M$ and $I_N\odot M_{\mathsf F}$ and their multiplication whiskerings. These domains fix which diagrams, rather than an unspecified list of axioms, are being checked.

\paragraph{Corner functors and paths.} The following notation suppresses only the already strong relative reassociators; composition of displayed functors is read right to left and juxtaposition is whiskering. Write $M=M_{\mathsf F}$. Set
\[
 D_2=(\mu_M N)(M\chi)(\chi M),\quad C_2=\chi(N\mu_M),\quad
 E_2=(M\mu_N)(\chi N)(N\chi),\quad A_2=\chi(\mu_N M).
\]
Thus the context comparison is $m:D_2\Rightarrow C_2$ and the content comparison is the isomorphism $n:A_2\Rightarrow E_2$.
For three context stages the source and target corner functors are
\[
 D_3=(\mu_M^{(3)}N)(MM\chi)(M\chi M)(\chi MM),\qquad
 C_3=\chi(N\mu_M^{(3)}).
\]
The first path applies $m$ to the first two stages, then to their composite and the third stage. The second applies $m$ to the last two stages after the first crossing, then to the first stage and their composite. Whiskering and strong multiplication associativity put both paths from $D_3$ to $C_3$. Equivalently their intermediate functors are $D_2(N\mu_M M)$ and $D_2(NM\mu_M)$, respectively, before the final $m$.
For three content stages the source is $\chi(\mu_N^{(3)}M)$ and the target is
\[
 (M\mu_N^{(3)})(\chi NN)(N\chi N)(NN\chi).
\]
The two paths expand $n$ first at the first pair or the last pair of content stages, then at the remaining pair, with the source multiplication associator.

For the mixed diagram all four corners are functors $NNMM\to MN$:
\[
 \begin{split}
 A&=\chi(\mu_N\mu_M),\qquad
 B=D_2(\mu_N MM),\\
 C&=E_2(NN\mu_M),\\
 D&=(\mu_M\mu_N)(M\chi N)(MN\chi)(\chi NM)(N\chi M).
 \end{split}
\]
The two paths are $B\xRightarrow{m}A\xRightarrow{n}C$ and
$B\xRightarrow{n^{[2]}}D\xRightarrow{m^{[2]}}C$. Here $n^{[2]}$ expands the two crossings of the composite content past the two context stages, using $n$ once at each crossing; $m^{[2]}$ contracts the two-stage context crossing once for each content layer, using $m$ at the inner and outer layer. Their other arrows are exactly whiskerings of the strong multiplications and their naturality. In the displayed $D$ the second and third crossings act on disjoint adjacent pairs and hence their alternative order is the ordinary functorial interchange. These formulas specify endpoints and pasted comparison generators, not merely a common domain word.

\begin{theorem}[Generating rewrite coherence]
The following diagrams commute on complete native data: naturality in native context/content arrows; compatibility with pure input/output actions and middle balance; both crossing-unit triangles; compatibility of the content-multiplication comparison with three successive content operations; compatibility of (1) with three successive context operations; and the mixed two-content/two-context diagram comparing content regrouping before versus after (1). Inserting unit generators into these diagrams gives the same forced-coordinate deletion maps. These comparisons induce the previously proved directed associator and its pentagon on $H_{\mathsf F}$.
\end{theorem}
\begin{proof}
Naturality and pure balance have already been checked at the coordinate functions. Their coefficient maps commute by tensor functoriality, and their whole grade maps commute by the original content naturality. In particular no route chooses a new target coupling independently of pushed source law.

For three content operations followed by one context, all routes flatten one finite deterministic occurrence tree to its addressed context leaves. Every route admits exactly the laws with the same leaf marginal constraints. Every raw map is that leaf-tuple bijection, every context coefficient map is the same labeled tensor coherence map, and the entire content grade map is the source substitution associator on every formal shape. This proves the three-content diagram, not merely equality on realized shape labels.

For one content followed by three context stages, take as common source the successively crossed, globally staged family and as common target the single crossing past the unary triple composite. The source constrains each new coordinate given the whole preceding stage tuples. The target tests each occurrence's complete local context path and its successive local prefix constraints. Each route applies (1) in one order, discarding only sibling-prefix observations that do not determine that coordinate's input type. At each deletion, convex conditional averaging proves admission. Both routes send the full raw triple tuple to the same occurrence-ordered triple tuple. Their context coefficient maps are the same permutation/reassociation of the actual first-, second- and third-stage labeled factors, and their full content grade maps are identity. Hence they are equal native arrows. At zero intermediate masses neither route specifies a conditional kernel; both retain the same full raw map there.

For two content operations followed by two context stages, the common raw data are pairs $(x_{ij},y_{ij})$ at the deterministic two-level occurrence leaves. One route first flattens the $(i,j)$ partition and applies (1); the other uses the invertible content comparison to retain blocks and applies the relevant local and outer versions of (1) before flattening. Both forget the same extra sibling first-stage observations, ultimately retaining the local $x_{ij}$ that fixes the second-stage input. Every generating arrow is admitted by the proved block-family equality or convex averaging. Both complete raw maps are the same leaf-tuple bijection. The context tensor arrows are the same labeled-factor map; the full grade routes are the same original $C\circ D$ coherence map. Thus the mixed diagram commutes.

Units insert only forced singleton coordinates. Their deletion commutes with flattening and discarding sibling observations, and actual unit coefficient arrows obey the monoidal unit equations. Thus the named crossing-unit and inserted-unit diagrams commute on the full data. These arguments establish equality of actual native morphisms because law actions are forced pushforwards and the common raw, context-coefficient, full grade and body fields have all been identified; no additional unspecified selector map remains.

Finally a mixed normal form is a context followed by content. Multiplying two such forms crosses the middle content/context pair and uses the two strong multiplications. Its raw law family is exactly the initializer-family mixed product: root marginal admitted, each continuation marginal conditional on the root raw atom admitted, arbitrary sibling coupling. Thus it is $H_{\mathsf F}$'s previously proved product. Comparing three mixed products uses precisely the two multiplication comparisons above and gives its one-way associator. For four mixed products, the listed three-stage and mixed diagrams, together with the strong context/content pentagons, give the same full raw flattening and actual labeled-factor maps; equivalently the already explicit mixed pentagon proof applies. This proves the directed composite coherence at the declared native dimension.\end{proof}

\paragraph{Why conditional coarsening is justified here.} In every directed step the finer prefix determines a coarser prefix that still fixes the new coordinate's input label and raw carrier. At positive coarser mass, its conditional law is a convex combination of the finer conditional laws in the same convex admitted family. This is the exact lemma used; discarding information that changes the coordinate's type would not satisfy it. It uses no measure-theoretic disintegration beyond finite sums and makes no claim about extra conditional choices on zero slices.

\section{Source comparison and remaining scope}
The singleton-family embedding $J_M$ of the original context module into $M_{\mathsf F}$ is strong on context multiplication and units: there is only one continuation at each unary context step, so fixed singleton outer and inner laws force the unique dependent product law. Thus singleton families are closed under unary dependent context accumulation. This is different from mixed composition across several content occurrences, which creates new sibling couplings. The entire content module $N$ is unchanged. The noninvertible source comparison first appears at the crossing: the original independent tuple law maps by identity raw/coefficient data into the all-coupling family. All source context coefficient arrows and all independent full content arrows remain retained.

Restrict each family context to its singleton distinguished source law. This is the homwise fully faithful source comparison already proved for the actual arity-preserving weighted-polynomial native arrows. The product member of the crossing family recovers the original independent crossing, with its original occurrence tensors, body equation and full grade. Product-law inclusions commute with the crossing comparisons above: both routes use the same original full independent law and the same actual tensor/grade maps. Their direction on mixed multiplication is
\[
 J_H(F;G)\longrightarrow J_H(F)\odot J_H(G),
\]
so this is the stated oplax/comonoidal comparison into a directed target, not a strong source inclusion into an ordinary bicategory.

The construction supplies a structured relative context/content explanation of the skew mixed product rather than only naming a row operation. It remains indexed by convex law families with common raw/coefficient maps and ordered position-preserving content maps. Arbitrary O1 content, separately chosen zero-prefix kernels, law-dependent raw maps, and full paid shared-world horizontal composition are not covered. The fixed-$P$ random-initializer feedback theorem is independently available, but its compatibility with this directed relative crossing or a paid-system extension still requires its own maps and equations; no automatic combination is asserted.
\end{document}
